\begin{definition}[The based comparison of the two B partners]
    \label{def:peps_torus_swept_patch_B_joint_loop}
    \lean{TNLean.PEPS.torusSweptPatchBJointLoop}
    \leanok
    \uses{def:peps_torus_swept_patch_endpoint_walks}
    Start at $B_1$, follow the actual partner connector to $B_2$, traverse its
    clockwise endpoint loop, and return along the reversed connector.
    Then traverse the clockwise endpoint loop at $B_1$.
    Thus, if $P$ is the connector transport, the reverse-ordered joint
    holonomy is $f_{B_1}(P^{-1}f_{B_2}P)$.
\end{definition}

\begin{theorem}[Actual partner transports before and after the gauge sweep]
    \label{thm:peps_torus_swept_patch_partner_transports}
    \lean{TNLean.PEPS.regularWalkHolonomy_torusSweptStringInitial_connectorA,
        TNLean.PEPS.regularWalkHolonomy_torusSweptStringInitial_connectorB,
        TNLean.PEPS.regularWalkHolonomy_torusSweptStringOperators_connectorA,
        TNLean.PEPS.regularWalkHolonomy_torusSweptStringOperators_connectorB}
    \leanok
    \uses{thm:peps_torus_swept_patch_connector_transport,
        thm:peps_torus_swept_string_transport,
        thm:peps_torus_swept_string_endpoint_holonomy}
    For the literal initial string assignment, the actual A and B connector
    transports are respectively $g^{-1}$ and $1$.
    After either of the two swept vertex gauges, they are respectively
    $g^{-1}$ and $g^{-1}$. These are identities for the same native walks
    at every translated position, including either seam.
\end{theorem}
\begin{proof}\leanok
    Substitute the literal directed bond values into the connector formulas.
    For either swept gauge, transport along a path changes by the gauge at
    its final base on the left and the inverse gauge at its initial base on
    the right. Both A bases remain outside the sweep. The B connector starts
    at the swept base and ends at its unchanged partner.
\end{proof}

\begin{theorem}[The same-connector gauge comparison has trivial joint flux]
    \label{thm:peps_torus_swept_patch_B_joint_holonomy}
    \lean{TNLean.PEPS.regularWalkHolonomy_torusSweptStringInitial_BJointLoop,
        TNLean.PEPS.regularWalkHolonomy_torusSweptStringOperators_transportedBPartner,
        TNLean.PEPS.regularWalkHolonomy_torusSweptStringOperators_BJointLoop}
    \leanok
    \uses{def:peps_torus_swept_patch_B_joint_loop,
        thm:peps_torus_swept_patch_partner_transports,
        thm:peps_regular_connector_holonomy,
        thm:peps_torus_swept_string_endpoint_holonomy}
    Initially the actual based B-partner comparison has joint holonomy $1$.
    After either gauge sweep, the transported inverse partner is
    $g h^{-1}g^{-1}$, and the joint holonomy is still
    \[
        (g h g^{-1})(g h^{-1}g^{-1})=1.
    \]
    This is an auxiliary statement about the literal swept gauge presentation
    in \cite[fluxon-braiding passage]{Schuch2010PEPS}. It does not identify
    physical reunion or string rerouting with this same connector, and is
    neither a counterexample to the source nor a physical braiding theorem.
\end{theorem}
\begin{proof}\leanok
    Transporting a loop with holonomy $f$ along a connector with transport $P$
    gives $P^{-1}fP$. Initially $P=1$ and the partner fluxes are $h,h^{-1}$.
    After the sweep $P=g^{-1}$, the first flux is $ghg^{-1}$ and the unchanged
    distant partner has flux $h^{-1}$. Their based product therefore cancels.
\end{proof}
